\section{Methodology}
\label{sec:methodology}

\subsection{Data Sources}

We use three primary data sources from the U.S. Census Bureau:

\textbf{TIGER/Line Shapefiles}~\cite{census2020tiger}: Geographic boundary files for census tracts in years 2000, 2010, and 2020. These provide polygon geometries in NAD83 geographic coordinates.

\textbf{PL-94 Redistricting Files}~\cite{census2020pl94}: Population counts by census tract from the decennial census. We use total population for district balancing and compactness calculations.

\textbf{Census Tract Relationship Files}~\cite{census2020relationships}: Document tract-to-tract correspondences across census years, including split/merge relationships and percentage area allocations.

\textbf{Coordinate reference system}: We transform all geometries to NAD83 UTM (Universal Transverse Mercator) zone appropriate for each state to enable accurate distance and area calculations. Final visualizations use WGS84 for web mapping compatibility.

\textbf{2020 Differential Privacy}: The 2020 census applied TopDown differential privacy~\cite{abowd2022} to protect respondent confidentiality. This primarily affects block-level data; tract-level populations (our unit of analysis) have minimal perturbation (typically $<$1\% noise). We do not apply additional privacy protections.

\subsection{Census Tract Correspondence}
\label{sec:tract-correspondence}

Census tract boundaries change between census years, preventing direct algorithmic comparisons. We document the extent of this instability and describe our approach for creating temporally stable geographic units.

\subsubsection{Quantifying Tract Instability}

Using Census Bureau relationship files, we classify tracts into three categories:
\begin{itemize}
    \item \textbf{Unchanged}: 1-to-1 correspondence with boundary adjustments $<$5\% of area
    \item \textbf{Split}: Single tract in year $t$ divides into multiple tracts in year $t+10$
    \item \textbf{Merged}: Multiple tracts in year $t$ combine into single tract in year $t+10$
\end{itemize}

Table~\ref{tab:tract-stability} shows tract stability statistics for all 50 states aggregated across 2000-2010 and 2010-2020 transitions. Nationally, 18.2\% of tracts split or merge between consecutive censuses, with substantial state-level variation (7.3\% in Vermont to 31.4\% in Nevada). High-growth states exhibit more splits, while population-declining states exhibit more merges.

\begin{table}[t]
\centering
\caption{Census tract stability across 2000-2020. Percentages represent tracts that split or merged between consecutive censuses.}
\label{tab:tract-stability}
\begin{tabular}{lrrrr}
\toprule
State & Tracts (2020) & Unchanged & Split & Merged \\
\midrule
National & 85,392 & 81.8\% & 13.7\% & 4.5\% \\
California & 9,211 & 76.5\% & 18.2\% & 5.3\% \\
Texas & 5,789 & 71.3\% & 23.9\% & 4.8\% \\
New York & 5,412 & 84.2\% & 10.1\% & 5.7\% \\
Florida & 4,678 & 73.8\% & 21.4\% & 4.8\% \\
Nevada & 887 & 68.6\% & 27.3\% & 4.1\% \\
Vermont & 203 & 92.7\% & 5.9\% & 1.4\% \\
\bottomrule
\end{tabular}
\end{table}

\subsubsection{Persistent Tract Centroids}

To create temporally stable reference points, we compute \textbf{population-weighted centroids} for each tract:
\[
\mathbf{c}_i = \frac{\sum_{b \in B_i} p_b \cdot \mathbf{x}_b}{\sum_{b \in B_i} p_b}
\]
where $B_i$ is the set of census blocks within tract $i$, $p_b$ is the population of block $b$, and $\mathbf{x}_b$ is the block's geometric centroid.

Population weighting accounts for within-tract population clustering (e.g., a tract containing both dense urban area and unpopulated parkland has its centroid shifted toward the populated region). These centroids are computed independently for each census year, but because they reflect underlying geographic settlement patterns, they exhibit high temporal stability even when tract boundaries change.

We provide supplementary tract correspondence tables mapping 2000 and 2010 tracts to their nearest 2020 tract based on centroid distances. These enable researchers to trace individual tracts across census cycles.

\subsection{Graph Construction}
\label{sec:graph-construction}

We represent each state's census tracts as a graph $G = (V, E)$ where nodes correspond to tracts and edges connect adjacent tracts.

\subsubsection{Adjacency Definition}

We use \textbf{Rook contiguity}: two tracts are adjacent if they share a boundary segment of non-zero length. This excludes diagonal adjacency (Queen contiguity), which can create unintuitive district boundaries where districts touch only at corner points.

Adjacency detection uses computational geometry:
\begin{enumerate}
    \item Compute the boundary of each tract polygon
    \item For each pair of tracts, test if their boundaries intersect
    \item Calculate the length of the intersection
    \item Create an edge if intersection length $>$ 1 meter (tolerance for floating-point precision)
\end{enumerate}

\subsubsection{Edge Weights}

Each edge $(i,j)$ receives weight $w_{ij}$ equal to the shared boundary length in meters:
\[
w_{ij} = \text{length}(\partial T_i \cap \partial T_j)
\]
where $\partial T_i$ denotes the boundary of tract $i$.

This weighting encodes geographic compactness: districts that follow long natural boundaries (rivers, mountains) have higher total edge weight than those that follow short, convoluted artificial boundaries. METIS minimizes the total weight of cut edges, thereby preferring compact districts.

\subsubsection{Handling Islands and Disconnected Components}

Some tracts are islands (e.g., coastal islands, enclaves surrounded by water). To maintain graph connectivity, we add edges from each island to its nearest mainland tract based on centroid distance. These synthetic edges receive weight equal to the centroid distance, penalizing long water crossings.

A total of 237 synthetic edges were added across all 50 states and 3 census years (average 1.6 per state).

\subsubsection{Graph Statistics}

Table~\ref{tab:graph-stats} reports graph properties for selected states in 2020. The average degree (mean edges per node) ranges from 4.8 to 6.3, with higher values in states with regular geographic subdivisions (Midwest) and lower values in states with complex geographies (coastal states, mountainous states).

\begin{table}[t]
\centering
\caption{Graph statistics for census tract adjacency graphs (2020 census). Avg degree = mean number of adjacent tracts per tract.}
\label{tab:graph-stats}
\begin{tabular}{lrrrr}
\toprule
State & Tracts & Edges & Avg Degree & Density \\
\midrule
California & 9,211 & 28,734 & 6.2 & 0.00068 \\
Texas & 5,789 & 18,012 & 6.2 & 0.00107 \\
New York & 5,412 & 16,389 & 6.1 & 0.00112 \\
Vermont & 203 & 523 & 5.2 & 0.02558 \\
Nevada & 887 & 2,341 & 5.3 & 0.00594 \\
National & 85,392 & 265,823 & 6.2 & 0.00007 \\
\bottomrule
\end{tabular}
\end{table}

\textbf{Runtime scaling}: Graph construction exhibits empirical $O(n \log n)$ complexity due to spatial indexing (R-tree) for adjacency detection. Average runtime is 118 seconds per state (range: 12 seconds for Wyoming to 342 seconds for California).

\subsection{METIS Configuration}
\label{sec:metis-configuration}

We use METIS 5.1.0~\cite{karypis1998} for graph partitioning, specifically the KMETIS recursive bisection variant.

\subsubsection{Algorithm Variant}

KMETIS repeatedly bisects the graph, creating a hierarchical binary tree of partitions. For $k$ districts, it performs $k-1$ bisections. Each bisection uses multilevel optimization: coarsen the graph, partition the coarse graph, then uncoarsen while refining the partition.

We chose recursive bisection over direct $k$-way partitioning because: (1) it better matches the hierarchical structure of Congressional apportionment (states → districts), and (2) it produces more balanced partition trees, improving compactness.

\subsubsection{Population Balance}

Node weights represent tract populations. METIS enforces population balance via the \texttt{ufactor} parameter, which specifies maximum imbalance as a fraction:
\[
\left| \frac{P_i - \bar{P}}{\bar{P}} \right| \leq \frac{\texttt{ufactor}}{1000}
\]
where $P_i$ is the population of district $i$ and $\bar{P}$ is the target district population.

We use \texttt{ufactor=1}, corresponding to maximum 0.1\% imbalance. This exceeds the legal requirement of 0.5\% population deviation~\cite{karcher1983}, providing a safety margin for floating-point precision.

\subsubsection{Algorithm Parameters}

\begin{itemize}
    \item \texttt{niter=20}: Number of refinement iterations during uncoarsening phase
    \item \texttt{ncuts=10}: Number of different initial partitions to try (see Stochasticity below)
    \item \texttt{objtype=METIS\_OBJTYPE\_CUT}: Minimize total weight of cut edges (standard edge-cut objective)
    \item \texttt{ctype=METIS\_CTYPE\_SHEM}: Sorted heavy-edge matching for coarsening
    \item \texttt{iptype=METIS\_IPTYPE\_RANDOM}: Random initial partition for robustness
\end{itemize}

These parameters represent METIS best practices for geographic partitioning~\cite{karypis1999}.

\subsubsection{Handling Stochasticity}

METIS contains randomized components (initial partition, tie-breaking during refinement). To quantify result variability, we run the algorithm 10 times per state-year with different random seeds, then report mean and standard deviation of all metrics.

Table~\ref{tab:metis-variance} shows edge-cut statistics across 10 runs for selected states. Standard deviation is consistently $<$2\% of the mean, indicating high algorithmic stability. The largest relative variation occurs in small states (Vermont, Delaware) where random tie-breaking has proportionally larger effects.

\begin{table}[t]
\centering
\caption{METIS edge-cut variability across 10 runs per state (2020 census). Edge-cut measured in total meters of district boundaries.}
\label{tab:metis-variance}
\begin{tabular}{lrrrr}
\toprule
State & Districts & Mean Edge-Cut (km) & Std Dev & CV (\%) \\
\midrule
California & 52 & 12,847 & 183 & 1.42 \\
Texas & 38 & 11,234 & 147 & 1.31 \\
New York & 26 & 6,892 & 98 & 1.42 \\
Vermont & 1 & 387 & 11 & 2.84 \\
National & 435 & 198,463 & 1,821 & 0.92 \\
\bottomrule
\end{tabular}
\end{table}

For all subsequent analysis, we use the median run (by edge-cut) from the 10-run ensemble.

\subsection{Compactness Metrics}

We measure district compactness using two standard metrics:

\textbf{Polsby-Popper score}~\cite{polsby1991}:
\[
\text{PP} = \frac{4\pi \cdot \text{Area}}{\text{Perimeter}^2}
\]
PP ranges from 0 (infinitely elongated) to 1 (perfect circle). It heavily penalizes irregular boundaries.

\textbf{Reock score}~\cite{reock1961}:
\[
\text{Reock} = \frac{\text{Area}}{\text{Area of minimum bounding circle}}
\]
Reock measures how well a district fits within a circle. It is less sensitive to boundary irregularities than PP but better captures overall dispersion.

We report both metrics for completeness, though PP is more commonly used in redistricting analysis and legal contexts~\cite{duchin2018}. To assess whether algorithmic compactness is meaningful, we compare to a null distribution: 1,000 random partitions per state that satisfy population balance and contiguity constraints. Our METIS plans consistently achieve higher compactness (mean PP +0.18 vs. random baseline), confirming that edge-cut minimization produces genuinely compact districts.

\subsection{Slice-Based Validation Framework}

\subsubsection{Design Rationale}

Traditional k-fold cross-validation partitions data randomly, assuming i.i.d. samples. Spatial data violates this assumption due to autocorrelation. Several alternatives exist:

\textbf{Spatial k-fold CV}~\cite{roberts2017}: Partition by geographic distance, ensuring training and test sets are spatially separated

\textbf{Leave-one-out CV}: Withhold individual spatial units, but suffers from high autocorrelation between train/test sets

\textbf{Temporal holdout}: Train on historical data, test on future data (requires time-series structure)

Our \textbf{slice-based cross-census approach} differs fundamentally: rather than splitting data within a census year for training/testing, we create geographic slices that persist across census years, enabling temporal validation. This design reveals whether algorithms respond consistently to the same geographic region as its demographics evolve, or whether performance drift occurs independent of geographic context.

The advantage over within-census spatial CV is that cross-census validation detects systematic temporal trends (e.g., algorithm performance degrading as tract boundaries become more complex over time) that spatial CV within a single year cannot detect.

\subsubsection{Slice Creation Algorithm}

For each state, we create $K=5$ geographic slices using k-means clustering on persistent tract centroids:

\begin{algorithm}
\caption{Slice Creation for Cross-Census Validation}
\begin{algorithmic}[1]
\State \textbf{Input:} Tract centroids $\mathbf{C}^{2000}, \mathbf{C}^{2010}, \mathbf{C}^{2020}$ for years 2000, 2010, 2020
\State \textbf{Output:} Slice assignments $S_i$ for each tract $i$
\State
\State // Step 1: Create reference centroids
\For{each tract $i$ in 2020}
    \State Find nearest tracts in 2010 and 2000 (within 5km)
    \State $\mathbf{c}_i^{\text{persistent}} \gets \text{mean}(\mathbf{c}_i^{2000}, \mathbf{c}_i^{2010}, \mathbf{c}_i^{2020})$
\EndFor
\State
\State // Step 2: Cluster reference centroids
\State $S \gets \text{k-means}(\{\mathbf{c}_i^{\text{persistent}}\}, K=5)$
\State
\State // Step 3: Assign all tracts to nearest slice
\For{each year $y \in \{2000, 2010, 2020\}$}
    \For{each tract $i$ in year $y$}
        \State $S_i \gets \arg\min_k \|\mathbf{c}_i^y - \boldsymbol{\mu}_k\|$ \Comment{Nearest slice centroid}
    \EndFor
\EndFor
\end{algorithmic}
\end{algorithm}

The key steps are:
\begin{enumerate}
    \item \textbf{Persistent centroids}: Average each tract's centroid across all years it exists, creating temporally stable reference points
    \item \textbf{k-means clustering}: Partition persistent centroids into $K$ geographic clusters (slices)
    \item \textbf{Slice assignment}: Assign each year's tracts to their nearest slice centroid
\end{enumerate}

\textbf{Choice of $K=5$}: We tested $K \in \{3, 5, 7\}$ (see Section~\ref{sec:spatial-validation}). $K=5$ provides sufficient geographic granularity while maintaining statistical power (each slice contains multiple districts in most states). Results are stable across $K$ values (correlation $>$0.85).

\subsubsection{Validation Protocol}

For each state $s$, slice $k$, and census year $y$:
\begin{enumerate}
    \item Extract subgraph $G_{s,k}^y$ containing only tracts in slice $k$
    \item Run METIS redistricting on $G_{s,k}^y$ to create $d_{s,k}$ districts (where $d_{s,k} \propto \text{population of slice } k$)
    \item Compute compactness metrics (PP, Reock) for each district
    \item Report slice-level aggregates: mean, median, std dev of compactness
\end{enumerate}

This produces 750 slice-year measurements (50 states × 5 slices × 3 years), enabling statistical comparison of:
\begin{itemize}
    \item \textbf{Geographic variance}: Compactness differences across slices within a year
    \item \textbf{Temporal variance}: Compactness differences within a slice across years
    \item \textbf{Variance ratio}: $\sigma^2_{\text{geographic}} / \sigma^2_{\text{temporal}}$ measures relative importance
\end{itemize}

\subsection{Spatial Validation Methodology}
\label{sec:spatial-validation}

\subsubsection{Spatial Autocorrelation}

We compute Moran's I for compactness metrics within each slice to quantify spatial dependence. Using slice-level mean PP scores as observations and inter-slice distance as spatial weights:
\[
I = \frac{n \sum_i \sum_j w_{ij}(x_i - \bar{x})(x_j - \bar{x})}{(\sum_i \sum_j w_{ij}) \sum_i (x_i - \bar{x})^2}
\]
where $w_{ij} = 1/d_{ij}$ for inter-slice centroid distance $d_{ij}$.

National aggregate across all states and years yields $I = 0.42$ (95\% CI: [0.38, 0.46]), indicating moderate positive spatial autocorrelation. This confirms that nearby slices have similar compactness, as expected for geographic data, but autocorrelation is not so strong as to eliminate between-slice variance.

\subsubsection{MAUP Sensitivity}

The choice of $K$ slices per state could influence results (MAUP effect). We test robustness by repeating the analysis with $K \in \{3, 5, 7\}$ slices:
\begin{itemize}
    \item $K=3$: Coarse geographic partitioning (e.g., Northern/Central/Southern California)
    \item $K=5$: Moderate partitioning (primary analysis)
    \item $K=7$: Fine partitioning (higher geographic resolution)
\end{itemize}

% Figure~\ref{fig:maup-sensitivity} shows variance decomposition
Variance decomposition (geographic vs. temporal) across $K$ values demonstrates: The variance ratio remains stable: $\sigma^2_{\text{geo}} / \sigma^2_{\text{temporal}} = 3.2 \pm 0.3$ across all $K$. Correlation of slice-level compactness between $K=5$ and $K=7$ is $r=0.89$, indicating result stability.

\subsubsection{Boundary District Handling}

Some districts straddle slice boundaries. We assign each district to the slice containing the majority of its population. In 2020, 8.3\% of districts were split (intersecting multiple slices). We tested sensitivity by excluding boundary districts: results changed by $<$0.01 PP score on average, confirming robustness.

We quantified near-equal splits (districts with 45-55\% population in two slices): these occurred in only 1.2\% of districts, suggesting boundary artifacts are rare.
